On 9 March 1497, Nicolaus Copernicus observed the occultation of Aldebaran by
the Moon from Bologna. In his work \textit{De revolutionibus orbium
c\oe lestium} (Book VI, Chapter 27) Copernicus described the event: ``I saw the
star touching the dark edge of the Moon and disappearing at the end of the 5th
hour of the night between the horns of the Moon, closer to the south horn by a
third of the Moon's diameter.''

Assuming that the occultation was observed on the local meridian, that at
maximum occultation Aldebaran was $0.32'$ above the southern edge of the Moon,
and that the apparent angular diameter of the Moon as seen from Bologna was
$31.5'$, solve the following tasks:

\begin{description}
\item[(a)] Find the latitude $\varphi_1$ of a place with the same longitude as
  Bologna, from which Aldebaran would have appeared to pass behind the centre
  of the Moon. (6 points)
\item[(b)] Find the duration of the occultation as seen from latitude
  $\varphi_1$ if Aldebaran appeared to pass along the diameter of the lunar
  disk. For simplicity, also assume that the Moon and the observer are moving
  linearly at constant speed, that the Moon's orbit is circular and that the
  declination of the Moon does not change during the occultation. (5 points)
\item[(c)] Find the topocentric angular velocity of the Moon against the
  background stars (for an observer on the surface of the Earth) during the
  occultation for an observer at latitude $\varphi_1$, in arcmin/hour, applying
  the same assumptions as in part (b). (7 points)
\item[(d)] Estimate the range of the Moon's topocentric angular velocities
  (against the background stars for an observer on the surface of the Earth) in
  arcmin/hour at latitude $\varphi_1$, assuming a circular orbit. Show how this
  result can be justified by expressing the relative velocity of the Moon and
  observer in terms of their velocity vectors. (7 points)
\end{description}

The declination of Aldebaran was $\delta_{\mathrm{A}} = 15.37^\circ$ in 1497
(due to precession), and the latitude of Bologna is
$\varphi_{\mathrm{B}} = 44.44^\circ$\,N.
